\documentclass[11pt]{article}
\usepackage[margin=1in]{geometry}
\usepackage{amsmath,amssymb,bm}
\usepackage[T1]{fontenc}
\usepackage[utf8]{inputenc}
\usepackage{lmodern}
\usepackage{booktabs}
\usepackage{longtable}
\usepackage{array}
\usepackage{enumitem}
\usepackage{hyperref}

\title{EKF Mathematical Formulation}
\author{sensor\_fusion documentation}
\date{\today}

\newcommand{\R}{\mathbb{R}}
\newcommand{\T}{^{T}}
\newcommand{\crossmat}[1]{\left[#1\right]_\times}
\newcommand{\dq}{\delta q}
\newcommand{\dt}{\Delta t}

\begin{document}
\maketitle
\tableofcontents
\newpage

\section{Scope}

This document describes the runtime EKF implemented by
\path{sensor_fusion/src/ekf/} and the included generated model files. The
filter estimates the NED-frame attitude with respect to the vehicle frame,
velocity, position, IMU biases, and the physical vehicle-to-body mount used
during inertial propagation.
It is the primary runtime inertial/GNSS filter behind \texttt{SensorFusion}.

\section{Frames and Mount Sources}

\begin{longtable}{>{\raggedright\arraybackslash}p{0.12\linewidth} >{\raggedright\arraybackslash}p{0.78\linewidth}}
\toprule
\textbf{Symbol} & \textbf{Meaning} \\
\midrule
\(n\) & Local NED frame. EKF position and velocity states are in this frame. \\
\(b\) & Raw IMU body frame. \\
\(v\) & Physical vehicle frame, FRD: forward, right, down. \\
\bottomrule
\end{longtable}

Rotations follow
\begin{equation}
  x_a=C_{ab}x_b,\qquad R(q_{ab})=C_{ab},\qquad
  R(q_1\otimes q_2)=R(q_1)R(q_2).
\end{equation}
The align filter supplies a physical vehicle-to-body mount. In the Rust field
layout this mount is stored in \path{q_bv0..q_bv3}, but its mathematical meaning
is \(q_{bv}\):
\begin{equation}
  C_{bv}=R(q_{bv}),\qquad x_b=C_{bv}x_v,\qquad
  C_{vb}=C_{bv}\T,\qquad x_v=C_{vb}x_b .
\end{equation}
The fusion facade exposes a single mount mode. In \texttt{MountMode::Auto},
Align seeds \(q_{bv}\) and the EKF continues estimating mount internally. In
\texttt{MountMode::Manual}, the caller supplies a fixed \(q_{bv}\) and EKF
mount states are frozen.

\section{Nominal and Error States}

The nominal state has 20 scalar components:
\begin{equation}
  x=
  \begin{bmatrix}
  q_{nv} & v_n\T & p_n\T & b_g\T & b_a\T & q_{bv}
\end{bmatrix}\T ,
\end{equation}
where \(q_{nv}\) rotates vehicle-frame vectors into NED, \(v_n\) is NED
velocity, \(p_n\) is local NED position, \(b_g\) and \(b_a\) are additive gyro
and accelerometer corrections in the raw IMU body frame, and \(q_{bv}\) is the
physical vehicle-to-body mount.

The 18-dimensional error state is
\begin{equation}
  \delta x =
  \begin{bmatrix}
  \delta\theta_v &
  \delta v_n &
  \delta p_n &
  \delta b_g &
  \delta b_a &
  \delta\psi_{bv}
  \end{bmatrix}\T .
\end{equation}
The Rust index order is:
\begin{center}
\begin{tabular}{ll}
\toprule
Indices & Error component \\
\midrule
0..2 & vehicle-attitude small angle \(\delta\theta_v\) \\
3..5 & velocity error \(\delta v_n\) \\
6..8 & position error \(\delta p_n\) \\
9..11 & gyro-bias error \(\delta b_g\) \\
12..14 & accelerometer-bias error \(\delta b_a\) \\
15..17 & vehicle-to-body mount error \(\delta\psi_{bv}\) \\
\bottomrule
\end{tabular}
\end{center}

\section{Nominal Propagation}

The EKF consumes one-sample raw body-frame IMU increments:
\begin{equation}
  u_k = \{\Delta\theta_b,\Delta v_b,\dt\}.
\end{equation}
Bias-corrected increments are
\begin{equation}
  \Delta\theta_b^c=\Delta\theta_b-b_g\dt,\qquad
  \Delta v_b^c=\Delta v_b-b_a\dt .
\end{equation}
The current mount rotates them into the vehicle frame:
\begin{equation}
  \Delta\theta_v^c=C_{vb}\Delta\theta_b^c,\qquad
  \Delta v_v^c=C_{vb}\Delta v_b^c,\qquad C_{vb}=C_{bv}\T .
\end{equation}
The generated nominal equations implement
\begin{align}
  q_{nv}^{+} &= \operatorname{normalize}\left(q_{nv}^{-}\otimes
  \dq(\Delta\theta_v^c)\right),\\
  v_n^{+} &= v_n^{-}+C_{nv}\Delta v_v^c+g_n\dt,\qquad
  C_{nv}=R(q_{nv}^{+}),\\
  p_n^{+} &= p_n^{-}+v_n^{+}\dt,\\
  b_g^{+}&=b_g^{-},\qquad b_a^{+}=b_a^{-},\qquad q_{bv}^{+}=q_{bv}^{-},
\end{align}
with \(g_n=[0,0,9.80665]\T\) in NED.

\section{Error-State Propagation}

The covariance prediction uses generated discrete matrices:
\begin{equation}
  \delta x_{k+1}=F_k\delta x_k+G_kw_k,
  \qquad
  P_{k+1}^{-}=F_kP_k^{+}F_k\T+G_kQ_kG_k\T .
\end{equation}
The 15-noise vector is
\begin{equation}
  w=
  \begin{bmatrix}
  n_g & n_a & n_{bg} & n_{ba} & n_m
  \end{bmatrix}\T .
\end{equation}
The diagonal covariance entries assembled at runtime are
\begin{equation}
  Q_k=\operatorname{diag}\left(
  \sigma_g^2\dt^2 I_3,\,
  \sigma_a^2\dt^2 I_3,\,
  \sigma_{bg}^2\dt I_3,\,
  \sigma_{ba}^2\dt I_3,\,
  \sigma_m^2\dt I_3
  \right).
\end{equation}
If mount states are frozen, the mount random-walk block is set to zero
and the mount covariance rows/columns are cleared.

\section{Measurement Updates}

Scalar observations use
\begin{align}
  r &= z-h(\hat x),\\
  S &= HPH\T+R,\\
  K &= PH\T S^{-1},\\
  \delta x &= Kr .
\end{align}
Covariance is updated with a scalar Joseph-form expression. The lateral and
vertical NHC rows may be fused as one sequential two-row batch so the second
row sees the covariance and accumulated correction from the first row.

\subsection{GNSS Position and Velocity}

GNSS position and velocity are local NED:
\begin{align}
  h_p(x)&=p_n,& H_p&=\begin{bmatrix}0&0&I_3&0&0&0\end{bmatrix},\\
  h_v(x)&=v_n,& H_v&=\begin{bmatrix}0&I_3&0&0&0&0\end{bmatrix}.
\end{align}
The generated GNSS rows have no direct mount Jacobian. The runtime has
optional mount update scales for GNSS rows; with the current generated \(H\)
these only affect mount updates induced through covariance cross-correlation.
The default public \path{fuse_gps} path passes zero scales.

\subsection{Zero Velocity}

The zero-velocity pseudo-measurement is a GNSS-velocity-shaped update:
\begin{equation}
  z_v=0,\qquad h_v(x)=v_n .
\end{equation}
It does not directly observe mount because it is expressed in NED, not in
vehicle coordinates.

\subsection{Vehicle Speed and NHC}

Predicted vehicle-frame velocity is
\begin{equation}
  v_v=C_{nv}\T v_n .
\end{equation}
The optional wheel/CAN speed measurement observes
\begin{equation}
  h_x(x)=e_x\T v_v .
\end{equation}
The non-holonomic constraints use zero lateral and vertical vehicle velocity:
\begin{equation}
  h_y(x)=e_y\T v_v\approx 0,\qquad
  h_z(x)=e_z\T v_v\approx 0 .
\end{equation}
Because mount is already part of propagation, these instantaneous vehicle-speed
rows have no direct mount Jacobian:
\begin{equation}
  \frac{\partial(e_i\T v_v)}{\partial\delta\psi_{bv}} = 0,\qquad i\in\{x,y,z\}.
\end{equation}
Mount is corrected through propagation-induced covariance with attitude and
velocity.

\subsection{Stationary Gravity}

Stationary gravity updates use vehicle-frame accelerometer samples and update
only \(x\) and \(y\) accelerometer channels:
\begin{equation}
  f_v \approx -C_{nv}\T g_n+\eta .
\end{equation}
The scalar rows are
\begin{equation}
  h_x=e_x\T(-C_{nv}\T g_n),\qquad
  h_y=e_y\T(-C_{nv}\T g_n).
\end{equation}
Before each stationary gravity scalar update, the runtime floors roll and pitch
covariance to avoid over-confident attitude locking. The mount Jacobian block
is zero for this pseudo-measurement.

\section{Injection and Covariance Reset}

After each update the error state is injected:
\begin{align}
  q_{nv}^{+} &= q_{nv}^{-}\otimes\dq(\delta\theta_v),\\
  v_n^{+} &= v_n^{-}+\delta v_n,\qquad
  p_n^{+}=p_n^{-}+\delta p_n,\\
  b_g^{+} &= b_g^{-}+\delta b_g,\qquad
  b_a^{+}=b_a^{-}+\delta b_a,\\
  q_{bv}^{+} &= \dq(\delta\psi_{bv})\otimes q_{bv}^{-}.
\end{align}
Both quaternions are normalized. The attitude and mount covariance
blocks are reset with
\begin{equation}
  G_r(\delta\theta)\simeq I-\frac{1}{2}\crossmat{\delta\theta}.
\end{equation}
The complete reset matrix is block diagonal with \(G_r\) on the attitude and mount
blocks and identity elsewhere:
\begin{equation}
  P^{+}\leftarrow G_{\mathrm{reset}}P^{+}G_{\mathrm{reset}}\T .
\end{equation}

\section{Initialization and Mount Modes}

At GNSS initialization the filter sets \(q_{nv}\) from GNSS heading or horizontal
velocity course, copies GNSS NED position and velocity into \(p_n,v_n\), sets
biases to zero, and initializes \(q_{bv}\) from Align in auto mode or from the
caller/reference source in manual mode. Covariance is diagonal with
configuration-controlled bias and mount variances.

The fusion facade controls the mount mode:
\begin{itemize}[leftmargin=1.5em]
\item \textbf{Auto}: Align seeds \(q_{bv}\) when the EKF initializes; the EKF
then estimates mount states internally.
\item \textbf{Manual}: a supplied or reference \(q_{bv}\) is used as a fixed
mount and EKF mount states are frozen.
\end{itemize}

\section{Known Modeling Boundaries}

\begin{itemize}[leftmargin=1.5em]
\item Mount affects inertial propagation directly by rotating raw IMU samples
into the vehicle frame. This is the primary mount observability path.
\item GNSS position, GNSS velocity, zero velocity, and stationary gravity do not
directly observe \(q_{bv}\).
\item Lateral and vertical NHC observations can couple attitude, velocity,
biases, and mount through covariance generated by propagation.
\item The physical mount used for plotting is the stored \(q_{bv}\) itself.
\end{itemize}

\end{document}
